\chapter{Variabili aleatorie multiple}

Molti esperimenti producono più numeri insieme: di una persona scelta a caso si possono misurare l'altezza, il peso e
l'età; da un canale di comunicazione esce un bit, ma interessa anche quale bit era entrato; una sorgente emette parole
di più bit. Questi numeri vanno descritti \textbf{insieme}, perché spesso sono legati fra loro: le persone alte
tendono a pesare di più, il bit ricevuto dipende da quello trasmesso. In questo capitolo introduciamo le coppie (e le
terne, e le $m$-ple) di variabili aleatorie, la loro \textbf{pmf congiunta}, le \textbf{marginali}, le
\textbf{condizionate} e l'\textbf{indipendenza}.

% ============================================================
\section{Coppie di variabili aleatorie}
% ============================================================

\dfn{Coppia di variabili aleatorie}{
  Una \term{coppia di variabili aleatorie} (o \term{variabile doppia}) associa a ogni risultato $\omega$ di un
  esperimento una coppia ordinata di numeri:
  \[ (X, Y) : \omega \in \Omega \;\longmapsto\; \left( X(\omega), Y(\omega) \right) \in \mathcal{X} \times \mathcal{Y}
     \subseteq \R^2 , \]
  dove $\mathcal{X}$ e $\mathcal{Y}$ sono gli alfabeti di $X$ e di $Y$.
}

In altre parole, il risultato di un esperimento non è più un unico valore $X(\omega) = x \in \mathcal{X}$, ma una
coppia ordinata $(X(\omega), Y(\omega)) = (x, y)$ che varia nel prodotto cartesiano $\mathcal{X} \times \mathcal{Y}$:
un punto del piano invece che un punto della retta.

\begin{figure}[H]
\centering
\begin{tikzpicture}
  \draw[mspigray, rounded corners=6pt] (-1.2,-1) rectangle (1.2,1);
  \node[font=\small, mspigray] at (-0.9,0.75) {$\Omega$};
  \fill (0.2,-0.1) circle (1.5pt) node[below, font=\small] {$\omega$};
  \begin{scope}[shift={(4,-1.2)}, x=0.55cm, y=0.55cm]
    \draw[-{Stealth}] (-0.3,0) -- (5.6,0) node[below, font=\small] {$x$};
    \draw[-{Stealth}] (0,-0.3) -- (0,4.4) node[left, font=\small] {$y$};
    \foreach \i in {1,...,5} \foreach \j in {1,...,4} \fill[mspigray!60] (\i,\j) circle (1.3pt);
    \fill[mspiB] (3,2) circle (2.4pt);
    \node[font=\small, mspiB!85!black, anchor=south west] at (3,2) {$(X(\omega), Y(\omega))$};
  \end{scope}
  \draw[-{Stealth}, mspiB, line width=0.8pt] (0.3,0) .. controls (2,0.6) and (4.5,0.8) .. (5.55,-0.1);
  \node[font=\scriptsize, align=center] at (6.6,-1.75) {i punti grigi sono $\mathcal{X} \times \mathcal{Y}$};
\end{tikzpicture}
\caption{Una coppia di variabili aleatorie manda ogni risultato in un punto del piano.}
\end{figure}

\begin{esempio}{Altezza, peso ed età}
Si consideri l'elenco di tutti i residenti in Italia a una certa data, con altezza ($X$), peso ($Y$) ed età ($Z$). Si
sceglie un residente a caso e se ne leggono l'altezza in centimetri, $X(\omega) \in \{30, 31, \ldots, 210\}$, il peso
in chilogrammi, $Y(\omega) \in \{0{,}5;\ 0{,}6;\ \ldots;\ 170\}$, e l'età in anni, $Z(\omega) \in \{0, 1, \ldots,
110\}$. Possibili coppie sono $(X, Y) \in \mathcal{X} \times \mathcal{Y}$, $(X, Z) \in \mathcal{X} \times
\mathcal{Z}$, $(Y, Z) \in \mathcal{Y} \times \mathcal{Z}$; e $(X, Y, Z) \in \mathcal{X} \times \mathcal{Y} \times
\mathcal{Z}$ è una \textbf{terna} di variabili aleatorie.
\end{esempio}

% ============================================================
\section{La pmf congiunta}
% ============================================================

Come per una variabile sola, la descrizione completa si ottiene dalle frequenze.
\begin{itemize}
  \item Si considera la coppia $(X, Y)$ (nell'esempio: altezza e peso di un residente scelto a caso).
  \item Si eseguono $n$ esperimenti (si sceglie $n$ volte a caso un nome dall'elenco).
  \item Si registrano i valori $(X(\tilde\omega_i), Y(\tilde\omega_i))$ (altezza e peso del residente scelto di volta
        in volta).
  \item Si indica con $n_{X=x,\, Y=y}$ il numero di volte in cui è uscito $X = x$ \emph{e} $Y = y$ (il residente aveva
        altezza $x$ e peso $y$): è il numero di occorrenze dell'evento $\{X = x\} \cap \{Y = y\}$.
\end{itemize}

\dfn{Pmf congiunta}{
  La \term{pmf congiunta} (o DF congiunta) delle due variabili aleatorie $X$ e $Y$ è
  \[ p_{X,Y}(x, y) = P\left(\{X = x\} \cap \{Y = y\}\right) = \lim_{n \to \infty} \frac{n_{X=x,\, Y=y}}{n} ,
     \qquad (x, y) \in \mathcal{X} \times \mathcal{Y} . \]
  Si scrive anche $P(X = x, Y = y)$, dove la virgola si legge «e».
}

La pmf congiunta è quindi una \textbf{tabella} di $\card{\mathcal{X}} \cdot \card{\mathcal{Y}}$ numeri, una casella per
ogni coppia di valori, che gode di certe proprietà.

\begin{esempio}{Una pmf congiunta su due bit}
Due variabili $X$ e $Y$ con alfabeto $\{0, 1\}$ hanno la pmf congiunta della tabella (la ritroveremo nel prossimo
capitolo):
\begin{center}
\begin{tikzpicture}
  \node at (0,0) {
    \renewcommand{\arraystretch}{1.5}
    \begin{tabular}{c|cc}
      $p_{X,Y}(x,y)$ & $y = 0$ & $y = 1$ \\ \hline
      $x = 0$ & $\frac13$ & $\frac29$ \\
      $x = 1$ & $\frac19$ & $\frac13$
    \end{tabular}};
  \begin{scope}[shift={(4.2,-1.1)}, x=1.6cm, y=1.6cm]
    \draw[-{Stealth}] (-0.3,0) -- (1.6,0) node[below, font=\small] {$x$};
    \draw[-{Stealth}] (0,-0.3) -- (0,1.6) node[left, font=\small] {$y$};
    \foreach \x/\y/\p/\l in {0/0/0.3333/\frac13, 0/1/0.2222/\frac29, 1/0/0.1111/\frac19, 1/1/0.3333/\frac13} {
      \fill[mspiA!45, draw=mspiA!85!black] (\x,\y) circle ({0.45*sqrt(\p)});
      \node[font=\scriptsize] at (\x,\y) {$\l$};
    }
    \node[font=\scriptsize, below] at (1,-0.05) {1}; \node[font=\scriptsize, left] at (-0.05,1) {1};
  \end{scope}
\end{tikzpicture}
\end{center}
A destra la stessa pmf disegnata come «bolle»: ogni punto $(x, y)$ ha un cerchio di area proporzionale alla sua
probabilità. I quattro numeri sono positivi e sommano a $\frac{3 + 2 + 1 + 3}{9} = 1$.
\end{esempio}

% ============================================================
\section{Proprietà della pmf congiunta}
% ============================================================

\teorema[congiunta]{Proprietà della pmf congiunta}{
  \[ p_{X,Y}(x, y) \geq 0 \qquad\text{e}\qquad \sum_{x \in \mathcal{X},\, y \in \mathcal{Y}} p_{X,Y}(x, y) = 1 . \]
}

\begin{dimostrazione}
\passo{Non negatività} È una probabilità.
\passo{La somma} In ogni prova la coppia assume esattamente un valore $(x, y)$: gli eventi elementari
$\{X = x, Y = y\}$ sono a due a due incompatibili (un evento $\{X = x, Y = y\}$ è incompatibile con ogni altro evento
$\{X = x', Y = y'\}$ con $(x', y') \neq (x, y)$) e la loro unione è $\Omega$. Per l'additività:
\[ 1 = P(\Omega) = P\left( \bigcup_{x \in \mathcal{X},\, y \in \mathcal{Y}} \{X = x, Y = y\} \right)
     = \sum_{x \in \mathcal{X},\, y \in \mathcal{Y}} P(\{X = x, Y = y\}) = \sum_{x,\, y} p_{X,Y}(x, y) . \]
\end{dimostrazione}

\subsection{Le pmf marginali}

Dalla pmf congiunta si ricavano le pmf delle singole variabili. Si noti che
\[ \bigcup_{x \in \mathcal{X}} \{X = x\} = \Omega , \qquad \bigcup_{y \in \mathcal{Y}} \{Y = y\} = \Omega : \]
qualche valore $X$ lo assume sempre, e anche $Y$.

\teorema[marginali]{Marginalizzazione della pmf congiunta}{
  \[ p_X(x) = \sum_{y \in \mathcal{Y}} p_{X,Y}(x, y) , \qquad\qquad p_Y(y) = \sum_{x \in \mathcal{X}} p_{X,Y}(x, y) . \]
  Le pmf $p_X$ e $p_Y$ ottenute così si chiamano \term{pmf marginali}.
}

\begin{dimostrazione}
\passo{Si «affetta» l'evento $\{X = x\}$} Intersecando con $\Omega = \bigcup_y \{Y = y\}$ e usando la proprietà
distributiva:
\[ \{X = x\} = \{X = x\} \cap \Omega = \{X = x\} \cap \bigcup_{y \in \mathcal{Y}} \{Y = y\}
   = \bigcup_{y \in \mathcal{Y}} \left( \{X = x\} \cap \{Y = y\} \right) . \]
\passo{Le fette sono incompatibili} Per $y \neq y'$ gli eventi $\{X = x\} \cap \{Y = y\}$ e $\{X = x\} \cap
\{Y = y'\}$ non possono verificarsi insieme ($Y$ non assume due valori). Quindi le probabilità si sommano:
\[ p_X(x) = P(\{X = x\}) = \sum_{y \in \mathcal{Y}} P(\{X = x\} \cap \{Y = y\}) = \sum_{y \in \mathcal{Y}} p_{X,Y}(x, y) . \]
\passo{Per $Y$} Lo stesso ragionamento, scambiando i ruoli, dà $p_Y(y) = \sum_x p_{X,Y}(x, y)$.
\end{dimostrazione}

Il nome viene da come si fanno i conti a mano: si sommano le righe e le colonne della tabella e si scrivono i totali
\textbf{a margine}.

\begin{figure}[H]
\centering
\renewcommand{\arraystretch}{1.6}
\begin{tabular}{c|cc|>{\columncolor{mspiB!12}}c}
  & $y = 0$ & $y = 1$ & $p_X(x)$ \\ \hline
  $x = 0$ & $\frac13$ & $\frac29$ & $\frac13 + \frac29 = \frac59$ \\
  $x = 1$ & $\frac19$ & $\frac13$ & $\frac19 + \frac13 = \frac49$ \\ \hline
  \rowcolor{mspiA!12} $p_Y(y)$ & $\frac13 + \frac19 = \frac49$ & $\frac29 + \frac13 = \frac59$ & 1
\end{tabular}
\caption{Le pmf marginali sono le somme delle righe (in arancione, la pmf di $X$) e delle colonne (in blu, la pmf di
$Y$), scritte ai margini della tabella della pmf congiunta.}
\end{figure}

Quindi caratterizzare congiuntamente la coppia $(X, Y)$ significa caratterizzare anche ciascuna variabile
marginalmente. Il viceversa \textbf{non} è vero in generale: le marginali non bastano a ricostruire la congiunta.

\begin{esempio}{Stesse marginali, congiunte diverse}
Due bit $X$ e $Y$, ognuno uguale a 0 o a 1 con probabilità $\frac12$. Le due pmf congiunte qui sotto hanno entrambe
queste marginali (si sommino righe e colonne), ma descrivono situazioni opposte: a sinistra i due bit non hanno nulla a
che fare l'uno con l'altro; a destra sono sempre uguali (per esempio, il secondo è una copia del primo).
\begin{center}
\renewcommand{\arraystretch}{1.5}
\begin{tabular}{c|cc}
  & $y = 0$ & $y = 1$ \\ \hline
  $x = 0$ & $\frac14$ & $\frac14$ \\
  $x = 1$ & $\frac14$ & $\frac14$
\end{tabular}
\hspace{2cm}
\begin{tabular}{c|cc}
  & $y = 0$ & $y = 1$ \\ \hline
  $x = 0$ & $\frac12$ & $0$ \\
  $x = 1$ & $0$ & $\frac12$
\end{tabular}
\end{center}
\end{esempio}

% ============================================================
\section{Variabili indipendenti}
% ============================================================

\dfn{Variabili aleatorie indipendenti}{
  Due variabili aleatorie $X$ e $Y$ sono \term{indipendenti} se (e solo se) gli eventi $\{X = x\}$ e $\{Y = y\}$ sono
  indipendenti per ogni coppia di valori $(x, y)$. Per due variabili indipendenti la pmf congiunta \textbf{si
  fattorizza} nel prodotto delle marginali:
  \[ p_{X,Y}(x, y) = P(\{X = x\} \cap \{Y = y\}) = P(\{X = x\})\, P(\{Y = y\}) = p_X(x)\, p_Y(y) . \]
}

Questo è l'unico caso in cui assegnare le due pmf marginali $p_X(x)$ e $p_Y(y)$ equivale ad assegnare la pmf
congiunta. Nell'ultimo esempio la tabella di sinistra è quella di due bit indipendenti (ogni casella vale
$\frac12 \cdot \frac12$); quella di destra no ($p_{X,Y}(0, 0) = \frac12 \neq \frac14$).

\begin{esempio}{Verificare l'indipendenza}
\begin{itemize}
  \item Due lanci di un dado onesto: $p_{X,Y}(i, j) = \frac{1}{36} = \frac16 \cdot \frac16 = p_X(i)\, p_Y(j)$ per ogni
        coppia: i due lanci sono indipendenti.
  \item La coppia dei due bit di prima, con marginali $p_X(0) = \frac59$ e $p_Y(0) = \frac49$:
        $p_{X,Y}(0, 0) = \frac13 = \frac{27}{81}$, mentre $p_X(0)\, p_Y(0) = \frac59 \cdot \frac49 = \frac{20}{81}$. Basta
        una sola casella che non si fattorizza: le due variabili \textbf{non} sono indipendenti.
\end{itemize}
\end{esempio}

Il concetto di pmf congiunta si generalizza a una $m$-pla di variabili aleatorie
$(X_1, X_2, \ldots, X_m) \in \mathcal{X}_1 \times \mathcal{X}_2 \times \cdots \times \mathcal{X}_m \subseteq \R^m$
mediante la pmf congiunta
\[ p_{X_1, X_2, \ldots, X_m}(x_1, x_2, \ldots, x_m) = P\left(\{X_1 = x_1\} \cap \{X_2 = x_2\} \cap \cdots \cap
   \{X_m = x_m\}\right) , \]
e il concetto di indipendenza si generalizza nel modo seguente: le $m$ variabili sono indipendenti se
\[ p_{X_1, \ldots, X_m}(x_1, \ldots, x_m) = \prod_{i=1}^{m} P(\{X_i = x_i\}) = \prod_{i=1}^{m} p_{X_i}(x_i) \]
per ogni scelta dei valori.

% ============================================================
\section{Le pmf condizionate}
% ============================================================

Si considerino le variabili aleatorie $X \in \mathcal{X}$ e $Y \in \mathcal{Y}$ con assegnata pmf congiunta
$p_{X,Y}(x, y)$. Applichiamo all'evento $\{X = x, Y = y\} = \{X = x\} \cap \{Y = y\}$ la legge della probabilità
composta:
\[ p_{X,Y}(x, y) = P(\{X = x\} \cap \{Y = y\}) = P(\{Y = y\} \mid \{X = x\})\, P(\{X = x\})
   = p_{Y \mid X}(y \mid x)\, p_X(x) . \]

\dfn{Pmf condizionate}{
  La \term{pmf condizionata di $Y$ dato $X$} è, per $p_X(x) > 0$,
  \[ p_{Y \mid X}(y \mid x) = P(Y = y \mid X = x) = \frac{p_{X,Y}(x, y)}{p_X(x)} , \]
  e allo stesso modo $p_{X \mid Y}(x \mid y) = p_{X,Y}(x, y) / p_Y(y)$ per $p_Y(y) > 0$.
}

Come la legge congiunta, la legge condizionata $p_{Y \mid X}(y \mid x)$ è una tabella di $\card{\mathcal{X}} \cdot
\card{\mathcal{Y}}$ numeri che soddisfa alcune proprietà. Siccome la congiunta si può scrivere condizionando in
entrambi i modi,
\[ p_{Y \mid X}(y \mid x)\, p_X(x) = p_{X,Y}(x, y) = p_{X \mid Y}(x \mid y)\, p_Y(y) , \]
si ottiene subito la versione per le pmf del teorema di Bayes.

\begin{formula}[legge di Bayes per le pmf]
\[ p_{X \mid Y}(x \mid y) = \frac{p_{Y \mid X}(y \mid x)\, p_X(x)}{p_Y(y)} \]
\end{formula}

\teorema[condpmf]{Proprietà delle pmf condizionate}{
  \begin{enumerate}[label=\alph*)]
    \item Fissato $y$ (con $p_Y(y) > 0$), al variare di $x$ la funzione $p_{X \mid Y}(x \mid y)$ è una pmf:
          $p_{X \mid Y}(x \mid y) \geq 0$ e $\sum_{x \in \mathcal{X}} p_{X \mid Y}(x \mid y) = 1$.
    \item Marginalizzazione con le condizionate:
          \[ p_X(x) = \sum_{y \in \mathcal{Y}} p_{X \mid Y}(x \mid y)\, p_Y(y) , \qquad
             p_Y(y) = \sum_{x \in \mathcal{X}} p_{Y \mid X}(y \mid x)\, p_X(x) . \]
    \item Se $X$ e $Y$ sono indipendenti: $p_{Y \mid X}(y \mid x) = p_Y(y)$ e $p_{X \mid Y}(x \mid y) = p_X(x)$.
  \end{enumerate}
}

\begin{dimostrazione}
\passo{a) È una pmf} È $\geq 0$ perché è una probabilità condizionata. Per la somma, gli eventi $\{X = x\}$ formano
una partizione di $\Omega$ e la probabilità condizionata è una probabilità (Teorema~\ref{teo:condprob}):
\[ \sum_{x \in \mathcal{X}} p_{X \mid Y}(x \mid y) = P\left( \bigcup_{x \in \mathcal{X}} \{X = x\} \;\Big|\; \{Y = y\}
   \right) = P(\Omega \mid \{Y = y\}) = 1 . \]
\passo{b) Marginalizzazione} Nella formula delle marginali si scrive la congiunta come condizionata per marginale:
$p_X(x) = \sum_y p_{X,Y}(x, y) = \sum_y p_{X \mid Y}(x \mid y)\, p_Y(y)$. Non è altro che la legge della probabilità
totale per l'evento $\{X = x\}$ rispetto alla partizione $\Omega = \bigcup_y \{Y = y\}$. Lo stesso per $p_Y$, rispetto
alla partizione $\Omega = \bigcup_x \{X = x\}$.
\passo{c) Indipendenza} Se $p_{X,Y}(x, y) = p_X(x)\, p_Y(y)$, allora $p_{Y \mid X}(y \mid x) = \frac{p_X(x)\,
p_Y(y)}{p_X(x)} = p_Y(y)$, e allo stesso modo per l'altra.
\end{dimostrazione}

\begin{figure}[H]
\centering
\begin{tikzpicture}
  \node (t) at (0,0) {
    \renewcommand{\arraystretch}{1.5}
    \begin{tabular}{c|c>{\columncolor{evid}}c}
      & $y = 0$ & $y = 1$ \\ \hline
      $x = 0$ & $\frac13$ & $\frac29$ \\
      $x = 1$ & $\frac19$ & $\frac13$ \\ \hline
      $p_Y(y)$ & $\frac49$ & $\frac59$
    \end{tabular}};
  \node (c) at (8.2,0) {
    \renewcommand{\arraystretch}{1.5}
    \begin{tabular}{c|c}
      & $p_{X \mid Y}(x \mid 1)$ \\ \hline
      $x = 0$ & $\frac{2/9}{5/9} = \frac25$ \\
      $x = 1$ & $\frac{1/3}{5/9} = \frac35$ \\ \hline
      somma & 1
    \end{tabular}};
  \draw[-{Stealth}, line width=0.8pt, mspiB] (t.east) -- node[above, font=\scriptsize, align=center, text width=2.6cm]
    {si prende la colonna $y = 1$ e la si divide per $p_Y(1) = \frac59$} (c.west);
\end{tikzpicture}
\caption{La pmf condizionata a $Y = 1$: si guarda solo la colonna $y = 1$ della congiunta e la si «rialza» dividendola
per il suo totale, così che sommi a 1. È la stessa operazione della pmf condizionata a un evento (Capitolo~6).}
\end{figure}

\begin{esempio}{Il canale binario con le pmf congiunte}
Riprendiamo il canale binario dei capitoli 3 e 4: $X$ è il bit trasmesso, con $p_X(0) = 0{,}6$ e $p_X(1) = 0{,}4$;
$Y$ è il bit ricevuto. Il canale è descritto dalla pmf condizionata $p_{Y \mid X}$: il bit arriva giusto con
probabilità $0{,}9$ e invertito con probabilità $0{,}1$. La congiunta si ottiene moltiplicando:
$p_{X,Y}(x, y) = p_{Y \mid X}(y \mid x)\, p_X(x)$.
\begin{center}
\renewcommand{\arraystretch}{1.4}
\begin{tabular}{c|cc|c}
  $p_{X,Y}(x,y)$ & $y = 0$ & $y = 1$ & $p_X(x)$ \\ \hline
  $x = 0$ & $0{,}9 \cdot 0{,}6 = 0{,}54$ & $0{,}1 \cdot 0{,}6 = 0{,}06$ & $0{,}6$ \\
  $x = 1$ & $0{,}1 \cdot 0{,}4 = 0{,}04$ & $0{,}9 \cdot 0{,}4 = 0{,}36$ & $0{,}4$ \\ \hline
  $p_Y(y)$ & $0{,}58$ & $0{,}42$ & 1
\end{tabular}
\end{center}
Tutto quello che avevamo calcolato è in questa tabella: la marginale $p_Y(1) = 0{,}42$ è la probabilità totale del
Capitolo~3, e la condizionata $p_{X \mid Y}(1 \mid 1) = \frac{0{,}36}{0{,}42} = \frac67$ è il risultato del teorema di
Bayes del Capitolo~4. La pmf congiunta è il modo più ordinato di tenere insieme tutte queste informazioni.
\end{esempio}

% ============================================================
\section{Più di due variabili: la regola della catena}
% ============================================================

Si consideri una terna di variabili aleatorie $(X, Y, Z)$, distribuite secondo $p_{X,Y,Z}(x, y, z)$,
$(x, y, z) \in \mathcal{X} \times \mathcal{Y} \times \mathcal{Z}$. Le pmf condizionate si definiscono come per le
coppie, e anche condizionando a più variabili insieme; per esempio
\[ p_{X \mid Y,Z}(x \mid y, z) = \frac{p_{X,Y \mid Z}(x, y \mid z)}{p_{Y \mid Z}(y \mid z)} \]
è la pmf di $X$ sapendo che $Y = y$ e $Z = z$.

\teorema[catena]{Regola della catena}{
  \[ p_{X,Y,Z}(x, y, z) = p_Z(z)\; p_{Y \mid Z}(y \mid z)\; p_{X \mid Y,Z}(x \mid y, z) . \]
  Ogni permutazione dei pedici e degli argomenti è ovviamente possibile, e la regola si estende a qualsiasi numero di
  variabili.
}

\begin{dimostrazione}
\passo{Prima applicazione della probabilità composta} Si considera l'evento $\{X = x, Y = y\}$ come un unico evento e
lo si condiziona a $\{Z = z\}$:
\[ P(X = x, Y = y, Z = z) = P(X = x, Y = y \mid Z = z)\, P(Z = z) . \]
\passo{Seconda applicazione, dentro il condizionamento} La probabilità condizionata a $\{Z = z\}$ è una probabilità
(Teorema~\ref{teo:condprob}), quindi anche per lei vale la probabilità composta:
\[ P(X = x, Y = y \mid Z = z) = P(X = x \mid Y = y, Z = z)\, P(Y = y \mid Z = z) . \]
\passo{Si mette insieme}
\[ P(X = x, Y = y, Z = z) = P(X = x \mid Y = y, Z = z)\, P(Y = y \mid Z = z)\, P(Z = z) . \]
\end{dimostrazione}

La regola della catena dice che una congiunta si può sempre costruire «un pezzo alla volta»: prima si sceglie $z$, poi
$y$ sapendo $z$, poi $x$ sapendo entrambi. È esattamente il diagramma ad albero, con un livello per ogni variabile.

\dfn{Indipendenza di più variabili}{
  Una terna di variabili aleatorie (e allo stesso modo un gruppo di $n > 3$ variabili) è \term{indipendente} se la pmf
  congiunta di un qualunque sottogruppo di esse, condizionata a un qualunque altro sottogruppo (costituito da
  variabili distinte dalle prime), è pari alla relativa pmf congiunta non condizionata.
}

In questo caso, chiaramente, la pmf congiunta si fattorizza nel prodotto delle pmf marginali: nella regola della
catena ogni condizionata diventa una marginale, e $p_{X,Y,Z}(x, y, z) = p_X(x)\, p_Y(y)\, p_Z(z)$. Attenzione: non
basta che le variabili siano indipendenti \emph{a coppie}, come mostra l'esempio che segue.

% ============================================================
\section{Esempio: tre bit da una sorgente binaria}
% ============================================================

Si consideri una sorgente binaria che emette tre bit, $(B_1, B_2, B_3)$, con $B_i \in \{0, 1\}$. Si assegnano due
possibili leggi congiunte, $p_{B_1,B_2,B_3}$ e $q_{B_1,B_2,B_3}$, riportate nella tabella ($0 < \alpha < 1$ è un
parametro). Ci chiediamo se le coppie $(B_1, B_2)$, $(B_1, B_3)$, $(B_2, B_3)$ e la terna $(B_1, B_2, B_3)$ siano o
meno indipendenti, secondo l'una e l'altra legge.

\begin{table}[H]
\centering
\renewcommand{\arraystretch}{1.3}
\begin{tabular}{@{}ccc@{}}
  \toprule
  $b_1 b_2 b_3$ & $p_{B_1,B_2,B_3}(b_1, b_2, b_3)$ & $q_{B_1,B_2,B_3}(b_1, b_2, b_3)$ \\
  \midrule
  000 & $(1-\alpha)^3$ & $(1-\alpha)^2$ \\
  001 & $\alpha(1-\alpha)^2$ & 0 \\
  010 & $\alpha(1-\alpha)^2$ & 0 \\
  011 & $\alpha^2(1-\alpha)$ & $\alpha(1-\alpha)$ \\
  100 & $\alpha(1-\alpha)^2$ & 0 \\
  101 & $\alpha^2(1-\alpha)$ & $\alpha(1-\alpha)$ \\
  110 & $\alpha^2(1-\alpha)$ & $\alpha^2$ \\
  111 & $\alpha^3$ & 0 \\
  \bottomrule
\end{tabular}
\caption{Le due leggi congiunte della sorgente di tre bit.}
\end{table}

\Needspace{18\baselineskip}
\subsection{La legge \emph{p}}

Cominciamo con il trovare la congiunta di tutte le possibili coppie secondo $p$. Per la marginalizzazione si somma
sul bit che non interessa; per esempio
\[ p_{B_1,B_2}(b_1, b_2) = p_{B_1,B_2,B_3}(b_1, b_2, 0) + p_{B_1,B_2,B_3}(b_1, b_2, 1) . \]
Nelle somme compare sempre il fattore $(1 - \alpha + \alpha) = 1$:
\begin{center}
\renewcommand{\arraystretch}{1.3}
\begin{tabular}{@{}cl@{}}
  \toprule
  $b_1 b_2$ & $p_{B_1,B_2}(b_1, b_2)$ \\
  \midrule
  00 & $(1-\alpha)^3 + \alpha(1-\alpha)^2 = (1-\alpha)^2 (1 - \alpha + \alpha) = (1-\alpha)^2$ \\
  01 & $\alpha(1-\alpha)^2 + \alpha^2(1-\alpha) = \alpha(1-\alpha)(1 - \alpha + \alpha) = \alpha(1-\alpha)$ \\
  10 & $\alpha(1-\alpha)^2 + \alpha^2(1-\alpha) = \alpha(1-\alpha)$ \\
  11 & $\alpha^2(1-\alpha) + \alpha^3 = \alpha^2(1 - \alpha + \alpha) = \alpha^2$ \\
  \bottomrule
\end{tabular}
\end{center}
Sommando sul bit 2 o sul bit 1 si trovano, allo stesso modo, $p_{B_1,B_3}$ e $p_{B_2,B_3}$: hanno esattamente gli
stessi quattro valori, $(1-\alpha)^2$, $\alpha(1-\alpha)$, $\alpha(1-\alpha)$, $\alpha^2$. Marginalizzando ancora,
$p_{B_1}(b_1) = p_{B_1,B_2}(b_1, 0) + p_{B_1,B_2}(b_1, 1)$:
\[
  p_{B_1}(0) = (1-\alpha)^2 + \alpha(1-\alpha) = (1-\alpha)(1 - \alpha + \alpha) = 1 - \alpha , \qquad
  p_{B_1}(1) = \alpha(1-\alpha) + \alpha^2 = \alpha ,
\]
e lo stesso vale per $B_2$ e $B_3$: ogni bit è 1 con probabilità $\alpha$. Ora si vede che la legge $p$ \textbf{si
fattorizza nel prodotto delle marginali}: ogni riga della tabella è un prodotto di tre fattori, $\alpha$ per ogni 1 e
$1 - \alpha$ per ogni 0 (per esempio $p(0, 1, 1) = (1-\alpha) \cdot \alpha \cdot \alpha$). Questo implica che anche le
congiunte a coppie si fattorizzano: secondo $p$ l'intera terna è statisticamente indipendente.

\subsection{La legge \emph{q}}

Procedendo allo stesso modo con $q$:
\begin{center}
\renewcommand{\arraystretch}{1.3}
\begin{tabular}{@{}cccc@{}}
  \toprule
  coppia di valori & $q_{B_1,B_2}$ & $q_{B_2,B_3}$ & $q_{B_1,B_3}$ \\
  \midrule
  00 & $(1-\alpha)^2$ & $(1-\alpha)^2$ & $(1-\alpha)^2$ \\
  01 & $\alpha(1-\alpha)$ & $\alpha(1-\alpha)$ & $\alpha(1-\alpha)$ \\
  10 & $\alpha(1-\alpha)$ & $\alpha^2$ & $\alpha^2$ \\
  11 & $\alpha^2$ & $\alpha(1-\alpha)$ & $\alpha(1-\alpha)$ \\
  \bottomrule
\end{tabular}
\end{center}
(Per esempio $q_{B_2,B_3}(1, 0) = q(0, 1, 0) + q(1, 1, 0) = 0 + \alpha^2$.) Le marginali sono
\[
  q_{B_1}(1) = q_{B_2}(1) = \alpha , \qquad
  q_{B_3}(0) = (1-\alpha)^2 + \alpha^2 , \qquad q_{B_3}(1) = 2\alpha(1-\alpha) .
\]
$B_1$ e $B_2$ sono indipendenti: $q_{B_1,B_2}$ è la stessa di prima, il prodotto delle marginali. Non lo sono invece
$B_1$ e $B_3$, né $B_2$ e $B_3$: per esempio $q_{B_1,B_3}(1, 0) = \alpha^2$, mentre
$q_{B_1}(1)\, q_{B_3}(0) = \alpha \left[ (1-\alpha)^2 + \alpha^2 \right]$, e i due numeri sono diversi (salvo il caso
speciale $\alpha = \frac12$, di cui diremo fra poco). Non lo è, quindi, nemmeno la terna.

Lo si vede bene con le condizionate. Sapendo che i primi due bit sono 1:
\[
  p_{B_3 \mid B_1 = 1, B_2 = 1}(1) = \frac{p_{B_1,B_2,B_3}(1, 1, 1)}{p_{B_1,B_2}(1, 1)} = \frac{\alpha^3}{\alpha^2} =
  \alpha = p_{B_3}(1) ,
\]
\[
  q_{B_3 \mid B_1 = 1, B_2 = 1}(1) = \frac{q_{B_1,B_2,B_3}(1, 1, 1)}{q_{B_1,B_2}(1, 1)} = \frac{0}{\alpha^2} = 0 \neq
  q_{B_3}(1) .
\]
Secondo $p$, conoscere i primi due bit non dice nulla sul terzo; secondo $q$, lo dice con certezza. Guardando la
tabella si capisce perché: nelle righe con probabilità non nulla, $B_3$ vale 1 esattamente quando $B_1$ e $B_2$ sono
diversi. Si vede facilmente che $B_3 = B_1 \oplus B_2$ è un \textbf{bit di parità}.

\begin{figure}[H]
\centering
\begin{tikzpicture}
  \node[draw, rounded corners=3pt, fill=mspiA!12, font=\small] (b1) at (0,0.6) {$B_1$};
  \node[draw, rounded corners=3pt, fill=mspiA!12, font=\small] (b2) at (0,-0.6) {$B_2$};
  \node[draw, circle, fill=mspiB!15, font=\large, minimum size=0.9cm, inner sep=0pt] (x) at (2.2,0) {$\oplus$};
  \node[draw, rounded corners=3pt, fill=mspiB!20, font=\small] (b3) at (4.2,0) {$B_3$};
  \draw[-{Stealth}] (b1) -- (x); \draw[-{Stealth}] (b2) -- (x); \draw[-{Stealth}] (x) -- (b3);
  \node[font=\scriptsize, align=center] at (0,1.35) {bit indipendenti,\\ 1 con probabilità $\alpha$};
  \node at (8.3,0) {
    \renewcommand{\arraystretch}{1.15}
    \begin{tabular}{cc|c}
      $b_1$ & $b_2$ & $b_1 \oplus b_2$ \\ \hline
      0 & 0 & 0 \\ 0 & 1 & 1 \\ 1 & 0 & 1 \\ 1 & 1 & 0
    \end{tabular}};
\end{tikzpicture}
\caption{La sorgente secondo la legge $q$: due bit indipendenti e un terzo bit calcolato come loro XOR (somma modulo
2), che vale 1 quando i due bit sono diversi. Il numero totale di uni nella parola è sempre pari: è il principio dei
codici di controllo di parità usati per rilevare gli errori di trasmissione.}
\end{figure}

\info{Indipendenti a coppie ma non tutti insieme}{
  Con $\alpha = \frac12$ succede una cosa curiosa. Tutte e tre le congiunte a coppie di $q$ valgono $\frac14$ in ogni
  casella, e le marginali valgono $\frac12$: \textbf{ogni coppia} di bit è indipendente. Eppure la terna non lo è
  affatto, perché il terzo bit è determinato dai primi due ($q(1,1,1) = 0 \neq \frac18$). Guardare i bit due alla
  volta non basta: per l'indipendenza di più variabili serve la fattorizzazione dell'intera pmf congiunta, come dice
  la definizione.
}

\gbox{In sintesi}{mspiA!80!black}{
\begin{itemize}
  \item Una \textbf{coppia di variabili aleatorie} manda ogni $\omega$ in un punto $(x, y)$ di
        $\mathcal{X} \times \mathcal{Y}$; la \textbf{pmf congiunta} $p_{X,Y}(x, y) = P(X = x, Y = y)$ è una tabella di
        numeri $\geq 0$ che somma a 1.
  \item \textbf{Marginali}: $p_X(x) = \sum_y p_{X,Y}(x, y)$, $p_Y(y) = \sum_x p_{X,Y}(x, y)$ (le somme «a margine»).
        La congiunta determina le marginali, non il viceversa.
  \item $X$ e $Y$ sono \textbf{indipendenti} se $p_{X,Y}(x, y) = p_X(x)\,p_Y(y)$ per ogni coppia: è l'unico caso in
        cui le marginali bastano.
  \item \textbf{Condizionate}: $p_{Y \mid X}(y \mid x) = p_{X,Y}(x, y)/p_X(x)$, una pmf in $y$ per ogni $x$; Bayes:
        $p_{X \mid Y} = p_{Y \mid X}\, p_X / p_Y$; marginalizzare con le condizionate è la probabilità totale.
  \item \textbf{Regola della catena}: $p_{X,Y,Z} = p_Z\, p_{Y \mid Z}\, p_{X \mid Y,Z}$. Più variabili sono
        indipendenti se la congiunta si fattorizza tutta; l'indipendenza a coppie non basta (bit di parità).
\end{itemize}
}
